\begin{table}[h]
\centering\small
\begin{threeparttable}
\caption{The substrates behind every number in this work, as median and the fifth to ninety-fifth percentile.}
\label{tab:si-applicability}
\begin{tabular}{lll}
\toprule
 & the evaluated test set & the comparison set \\
\midrule
molecular weight & 333.9 ($138.64$, $662.77$) & 342.35 ($138.14$, $613.1$) \\
calculated logP & 3.01 ($-0.53$, $6.42$) & 2.9 ($0.12$, $6.04$) \\
heavy atoms & 23.5 ($9$, $46.55$) & 24 ($10$, $42$) \\
rings & 3 ($0$, $6$) & 3 ($0.5$, $6$) \\
rotatable bonds & 4 ($0$, $14$) & 4 ($0$, $11.5$) \\
\midrule
substrates & 1170 & 291 \\
Bemis--Murcko scaffolds & 579 & 207 \\
scaffolds on one substrate & 399 & 173 \\
acyclic substrates & 68 & 15 \\
\bottomrule
\end{tabular}
\begin{tablenotes}[flushleft]\footnotesize
\item The corpus is assembled from four drug-centric sources and the distribution is what that produces; the manuscript motivates the problem with environmental chemicals as well, and this table is what a reader should hold that motivation against. It describes the population and does not establish an applicability domain, which would need performance measured against these axes rather than the axes alone.
\end{tablenotes}
\end{threeparttable}
\end{table}
